\section{Conclusion}
\label{sec:conclusion}

This paper presented a performance evaluation of a Go-native Sparse
Merkle Tree with Poseidon2 hashing for the Basis Network enterprise
zkEVM L2 state database. Our key findings are:

\begin{enumerate}
  \item \textbf{Go achieves 10--14$\times$ speedup over TypeScript}
    for all SMT operations, driven by gnark-crypto's
    assembly-optimized BN254 field arithmetic. Poseidon2 hashing at
    4.46~$\mu$s per 2-to-1 compression enables 224,000 hashes per
    second.

  \item \textbf{Block-level state root computation meets the 50~ms
    target} for up to 250 transactions per block on a 10,000-entry
    tree at depth 32, with batch latency of 46.05~ms (mean) and
    48.92~ms (P95).

  \item \textbf{Depth sensitivity is the critical constraint.}
    At production EVM depths (160--256), the 50~ms budget is only
    achievable for 20--50 transactions per block without compact
    SMT optimizations or batch parallelism.

  \item \textbf{Memory efficiency is excellent} at 2.9~KB per
    account (29.3~MB for 10,000 entries), with linear scaling
    enabling in-memory operation for typical enterprise deployments.

  \item \textbf{Hash function alignment is architecturally critical.}
    Poseidon2 (gnark-crypto) and original Poseidon (circomlibjs/arbo)
    produce incompatible hashes. The state database hash function
    must match the prover's circuit library.
\end{enumerate}

The hypothesis is confirmed with the qualification that production
deployments at full EVM depths require compact SMT implementations
with path compression and batch parallelism---standard optimizations
in production zkEVM systems (Polygon, Scroll, zkSync).

\subsection{Recommendations for Implementation}

Based on these findings, we recommend the following for the Basis
Network zkEVM L2 state database:

\begin{enumerate}
  \item Implement a compact SMT (Jellyfish-style) with path
    compression to achieve $O(\log n)$ effective depth.
  \item Select gnark-crypto with Poseidon2 if the prover uses
    gnark/halo2; select vocdoni/arbo with original Poseidon if
    the prover uses circom/snarkjs.
  \item Use LevelDB or Pebble for persistent storage with an
    in-memory cache for hot accounts.
  \item Implement batch parallelism for blocks exceeding 100
    transactions, exploiting disjoint subtree independence.
  \item Target 100--250 transactions per block as the operational
    sweet spot balancing throughput and state root latency.
\end{enumerate}

\subsection{Future Work}

Three directions emerge from this study:

\begin{itemize}
  \item \textbf{Compact SMT benchmarking:} Quantify the speedup from
    path compression at depths 160 and 256 with realistic EVM key
    distributions.
  \item \textbf{Persistent storage integration:} Measure the
    performance impact of LevelDB/Pebble-backed storage vs.\
    in-memory operation.
  \item \textbf{Witness generation coupling:} Profile the end-to-end
    latency from state update through witness extraction to ZK
    circuit input preparation.
\end{itemize}

These investigations will be pursued in subsequent research units
as the Basis Network zkEVM L2 pipeline progresses toward production
deployment.
